\section{Handshake Specification of the Source-Code Versions}
\label{app:handshake-spec}

This appendix records the field-level specification of the handshake
implemented by the two source-code versions, summarised in
Section~\ref{subsubsec:variant-spec}. It is reference material: the
design rationale, security requirements, and implementation status are
stated in Sections~\ref{subsec:objectives}
to~\ref{subsec:variant-design}.

The responder holds four kinds of key material: a long-term ML-DSA-87
identity pair; a medium-term X25519 ratchet pair; a signed last-resort
KEM pair; and a replenished collection of signed one-time KEM pairs.
The ratchet key does not enter the KEM-only bootstrap secret, but it
seeds the Double Ratchet and contributes to the post-bootstrap message
chain. The initiator holds its own long-term ML-DSA-87 identity pair
and creates a fresh X25519 base key for each run, which identifies the
session and is included in the signed transcript. Every public key
begins with a one-byte algorithm identifier: \texttt{0x09} for the
ML-DSA-87 identity keys, and, for the KEM keys, \texttt{0x0A}
(ML-KEM-1024) in Version 1 and \texttt{0x0B} (HQC-256) in Version 2.

The responder's published bundle contains its identity verification
key, the ratchet public key and signature, the last-resort KEM public
key and signature, an optional one-time KEM public key and signature,
and the numeric identifiers of those signed keys; each signature
covers the encoded key but not its numeric identifier or assigned
bundle role. The initial message extends libsignal's
\texttt{PreKeySignalMessage} with three fields: a one-time KEM pre-key
identifier, a one-time KEM ciphertext, and the initiator's transcript
signature. Existing fields carry the initiator's identity and base
keys, the last-resort ciphertext and key identifiers, and the
encrypted first message.

Encapsulation to the last-resort key yields \(ss_1\) and, when a
one-time key is used, encapsulation to it yields \(ss_2\). The
initiator forms the secret input
\(\mathtt{0xFF}^{32} \parallel ss_1 \parallel ss_2\), omitting
\(ss_2\) in last-resort-only mode; the leading constant preserves
libsignal's discriminator convention. HKDF-SHA-256 derives 64 bytes of
bootstrap root and chain material with no salt; the \texttt{info}
value is \texttt{PQXDH\_MLKEM1024\_MLDSA87\_SHA-256} in Version 1 and
\texttt{PQXDH\_HQC256\_MLDSA87\_SHA-256} in Version 2. The signed
transcript is
\[
t = \mathrm{lp}(L_v) \parallel \mathrm{lp}(IK_A) \parallel
\mathrm{lp}(ct_1) \parallel \mathrm{lp}(ct_2) \parallel
\mathrm{lp}(IK_B) \parallel \mathrm{lp}(RK_B) \parallel
\mathrm{lp}(BK_A),
\]
where \(\mathrm{lp}(\cdot)\) prefixes each value with its four-byte
big-endian length; \(L_v\) is
\texttt{PQXDH\_MLKEM1024\_MLDSA87\_transcript} or
\texttt{PQXDH\_HQC256\_MLDSA87\_transcript}; \(ct_2\) is the empty
string in last-resort-only mode; \(RK_B\) is the responder's ratchet
public key; and \(BK_A\) is the initiator's base key. Every ML-DSA-87
operation uses the context string \texttt{Signal\_PQXDH\_MLDSA87}.

\section{Extended Requirements and Acceptance Criteria}
\label{app:extended-requirements}

This appendix expands the concise requirements of Section~\ref{subsec:objectives}. It provides supporting specification details and audit material. The principal claim boundaries and outcomes remain in Chapter~\ref{sec:design}.

\subsection{Functional and Test Acceptance Criteria}

Each source-code version must execute the complete PQXDH-derived handshake, from bundle generation through signature verification, encapsulation, decapsulation, transcript authentication, and secret derivation. The design must state how identity verification keys are bound to users and must authenticate the algorithm-suite, key-role, and key identifiers needed to interpret the transcript. Each public key and transcript field must have one clearly defined byte representation. The encoding must also identify the type and purpose of each value so that the same bytes cannot be interpreted in different ways. The initial encrypted message must be authenticated before it is accepted. Transient shared secrets and consumed one-time private pre-keys must be erased as soon as they are no longer required.

Automated tests must verify successful operation in both implemented pre-key modes. Negative tests must separately alter every signature, signed public key, identity or algorithm identifier, KEM ciphertext, authenticated transcript field, and initial encrypted message, and must reject malformed encodings and inputs from the other source-code version. Tests must also resend an unchanged copy of the initial message to check that it cannot create a second live session or cause the same plaintext to be delivered twice. If the handshake itself does not reject this replay, the design must identify the surrounding replay-protection mechanism responsible for doing so.

The artefact's handling of post-quantum pre-keys differs from deployed PQXDH, which selects one signed KEM pre-key per run: a one-time post-quantum pre-key when one remains, or otherwise the signed last-resort pre-key. Each artefact version instead always encapsulates to the signed last-resort KEM pre-key and, when available, additionally to a one-time KEM pre-key, so the two implemented modes are last-resort-only and last-resort-plus-one-time rather than deployed PQXDH's either/or selection. The optional curve one-time pre-key of deployed PQXDH has no counterpart in these no-X25519 handshake KDFs.

\subsection{Security and Operational Acceptance Criteria}

Forward-secrecy claims depend on the pre-key mode, on authenticated pre-key role metadata, and on timely key erasure. In last-resort-only mode, later compromise of the retained last-resort decapsulation key can expose a recorded handshake that used it. In last-resort-plus-one-time mode, the one-time private key and its transient shared secret must be deleted after successful use. Under the KEM and KDF assumptions, later compromise of the retained last-resort key alone must then not suffice to recover the initial handshake secret. The stronger guarantee also requires both parties to authenticate that the additional key is genuinely one-time, a condition whose implementation status Section~\ref{subsubsec:design-deniability} records. Compromise of a long-term identity signing key after completion must not by itself reveal a previously derived handshake secret, although it may enable later impersonation.

The handshake must satisfy four further targets, acceptance criteria rather than properties inherited from the PQXDH proof. Session independence: compromise of one completed handshake's secrets must not aid an adversary against any other; each run must draw fresh encapsulation randomness and reuse no other run's shared secrets. Key-compromise impersonation resistance: compromising one party's identity signing key must not enable impersonation of a different, uncompromised party to the compromised party. Identity misbinding resistance: a completed handshake must commit both parties to a consistent view of the two identity keys and the exchanged key material, so neither can attribute the session to an unintended peer. Replay and key reuse: an exact replay of a recorded initial message must not establish a second live session or deliver plaintext twice, and may be handled by a defined mechanism in the surrounding session layer (Section~\ref{subsubsec:variant-spec}); a one-time pre-key must serve at most one accepted handshake, while reuse of the signed last-resort pre-key across sessions is inherent to its function and must be restricted to that role.

Three operational requirements follow. All key generation and encapsulation randomness must come from a cryptographically secure source, including the seeds of deterministic entry points. Transient shared secrets, the assembled KDF input, and the private halves of consumed one-time pre-keys must be promptly erased once no longer needed. A malicious server that withholds one-time pre-keys forces last-resort-only mode and its weaker forward-secrecy profile. The initiator cannot detect the withholding, a limitation inherited from deployed PQXDH \cite{kret_pqxdh_nodate}, so the downgrade is accepted and explicitly modelled in Section~\ref{subsec:threat-model}.

These requirements rest on stated assumptions about the surrounding primitives, consistent with the PQXDH specification and its formal analyses \cite{kret_pqxdh_nodate, bhargavan_formal_2024}: an authenticated-encryption construction for the initial message providing IND-CPA confidentiality and INT-CTXT ciphertext integrity against a quantum-capable adversary; HKDF modelled as a secure key derivation function on the specified inputs; pairwise-disjoint encodings for all key types; validation of every public input prior to use; a cryptographically secure random source at both endpoints; and reliable deletion, via zeroisation, of expired secret material. The artefact adopts libsignal's AES-256-CBC encrypt-then-MAC construction, HMAC-SHA-256 truncated to an eight-byte tag and verified before decryption, rather than a standalone category-5 AEAD. Whether this inherited construction, with its 64-bit tag bound, meets the abstract post-quantum integrity assumption requires separate justification and is not a demonstrated property. Section~\ref{subsubsec:variant-spec} records which of these the artefact demonstrates and which remain obligations for the surrounding application.

\subsection{Role-Level Signing Requirements}

What must be signed is specified by role. The responder's offline signatures must bind each published pre-key to the responder's identity and to the role, algorithm suite, and key identifiers necessary for interpretation; the initiator's initial-message signature must bind both identities, the selected pre-key identifiers and public keys, the KEM ciphertext or ciphertexts, and the complete canonical transcript. The division arises from asynchronous operation: the responder is offline when the initiator encapsulates, so it cannot sign a subsequent ciphertext, and ciphertext binding falls to the initiator.

\section{KEM Binding Mechanisms in Detail}
\label{app:binding-detail}

This appendix records the mechanism-level evidence supporting the binding verdicts in Section~\ref{subsubsec:binding-check}. Structural observations are not substituted for the published proof status.

ML-KEM-1024 passes the primitive-level binding gate. During encapsulation it computes \((K,r)=G(m \parallel H(ek))\), where \(H\) is SHA3-256, \(G\) is SHA3-512, \(H(ek)\) hashes the complete encapsulation key, \(K\) is the shared secret, and \(r\) is the randomness used to produce the ciphertext. During decapsulation it decrypts the message, recomputes these values, and deterministically re-encrypts it; if the recomputed ciphertext does not match, implicit rejection returns a replacement secret without revealing the failure \cite{national_institute_of_standards_and_technology_us_module-lattice-based_2024-1}. These structural features alone are not treated as the proof; the selection relies on the published PQXDH-specific result of Fiedler and G\"unther \cite{fiedler_security_2024}.

HQC requires a revision-specific conclusion. The published binding analysis applies to the round-four submission and finds that construction deficient under several binding notions, including in the honest-key setting \cite{noauthor_binding_nodate}. Version 2 instead implements v5.0.0, in which encapsulation computes \((K,\theta)=G(H(ek_{\mathrm{KEM}}) \parallel m \parallel salt)\). Here, \(H\) is a domain-separated SHA3-256 hash of the complete encapsulation key, while \(G\) is a domain-separated SHA3-512 hash. The first 32 bytes of \(G\)'s output form the shared secret \(K\), and the remaining 32 bytes form the encryption randomness \(\theta\); \(m\) is a random 32-byte message, and the random 16-byte salt is carried in the ciphertext. During decapsulation, HQC decrypts \(m\), recomputes the same values, and re-encrypts it. If the recomputed ciphertext does not match the received one, it returns a replacement secret derived from the received ciphertext instead of revealing the failure. The encapsulation key and salt therefore enter the successful shared-secret derivation in this revision, changes documented by the updated specification and NIST's FIPS 207 conversion presentation \cite{aguilar-melchor_hqc_nodate,csrcdatamodelspresentationspresentersjsonobject_csrc_2025}; direct inspection confirms that the pinned Rust source follows this construction. The published binding paper analyses the earlier round-four scheme, not v5.0.0, so these changes are supporting evidence rather than proof that Version 2 satisfies the exact notion PQXDH requires.

\section{Operational Modes and Secret Erasure}
\label{app:spec-operational}

This appendix records implementation-level operational evidence supporting Sections~\ref{subsubsec:req-functional} and~\ref{subsubsec:variant-spec}.

The two operational modes are last-resort-only and last-resort-plus-one-time. When the store has no one-time key, or a malicious server withholds it, the handshake falls back to the weaker profile of Appendix~\ref{app:extended-requirements}, and the stronger one-time claim remains conditional while the signatures do not authenticate a key's role or identifier. Replay protection is delegated to the surrounding session state: with state retained, the Double Ratchet's message-key accounting is expected to reject the repeated payload, but no end-to-end test shows that a replayed complete initial message is rejected, and state loss or rollback can allow a last-resort-only initial message to be accepted again. Replay resistance remains an unverified application-layer obligation.

Secret erasure has clear requirements but incomplete implementation evidence. Each side should erase \(ss_1\), \(ss_2\), and the assembled HKDF input immediately after derivation; the responder should delete a one-time private key after its first successful decapsulation; and the initiator should discard pending ciphertexts and the transcript signature once the responder's first reply is processed. Inspection of the pinned HQC crate shows zeroisation of selected seeds, decrypted messages, and derivation intermediates, but not of every adapter or protocol-layer copy. One-time-key consumption is delegated by libsignal to a store hook that the bundled in-memory store does not act on. Protocol-layer zeroisation of the shared secrets and HKDF input is not demonstrated. These are recorded as unmet requirements rather than properties of the implementation.

\section{Size Accounting Model}
\label{app:size-accounting}

This appendix defines the component sums and variables used by Section~\ref{subsec:theoretical-comparison}; it does not redefine them as complete wire or service encodings.

The accounting follows the artefact's encodings. Every public key and KEM ciphertext carries a one-byte algorithm identifier, while ML-DSA-87 and XEdDSA signatures carry none; an ML-DSA-87 identity verification key therefore occupies 2593 bytes, 2592 from FIPS 204 plus the \texttt{0x09} type byte. A bundle cryptographic-component sum includes the identity key, the ratchet or signed pre-key and its signature, every KEM pre-key and its signature, and the baseline's optional one-time curve key; an initial-message cryptographic-component sum includes the initiator's identity and base or ephemeral keys, the KEM ciphertext or ciphertexts, and, for Versions 1 and 2, the transcript signature. Numeric identifiers, Protocol Buffers framing, and the encrypted first message are excluded; these fields are not assumed equal between configurations, which is why Chapter~\ref{sec:methodology} measures complete initial-message and reply protocol objects separately and reports a serialised bundle-component total rather than a complete bundle encoding.

Modelled server-held public cryptographic payload is reported per device. For Versions 1 and 2 it is parameterised by \(N\), the number of stored one-time KEM pre-keys; for the baseline, the one-time KEM and curve-key inventories are represented separately by \(K\) and \(C\). A single unqualified per-user value would otherwise conceal the largest variable components. Table~\ref{tab:spec-sizes} combines the KEM and signature sizes established in Sections~\ref{subsubsec:kem-selection} and~\ref{subsubsec:sig-selection}. The baseline uses 33-byte encoded X25519 public keys and 64-byte XEdDSA signatures, and follows PQXDH's either/or selection of a one-time or last-resort Kyber pre-key \cite{kret_pqxdh_nodate}. Its per-run KEM material is consequently one public key and one ciphertext in either mode.